% =====================================================================
%  VIII. Conclusion
%
%  Khong duoc ket bang "ket qua hua hen". Ket bang cai ta do duoc va cai
%  ta khong do duoc. Cau cuoi cung phai la mot cau kiem chung duoc.
% =====================================================================

\revnote{rewrite}{Rewritten: the advantage claim is replaced by the boundary the regime map draws (AE-1, R1-4, R3-2).}
\begin{revbody}
We set out to locate, rather than to demonstrate, the value of a NISQ-feasible
quantum kernel for network intrusion detection. Varying one factor at a time
over $110$ controlled comparisons, we find $21$ verdicts favouring the quantum
kernel, $21$ favouring a classical baseline, and $68$ inconclusive when each
is corrected within its baseline family; $17$, $15$ and $78$ under one
Benjamini--Hochberg correction over the map; and no surviving verdict at all
under Holm over all $110$ cells. Inside that balance there is a structure that
a single-operating-point benchmark cannot see, and it is stated below with the
protocol each part of it comes from.

The structure has three parts. The ordering between \QSVM{} and the tree
ensembles and the linear and polynomial kernels reverses between $N=2000$ and
$N=5000$ under the natural class prior, in both tuning arms, on a shared
four-component embedding; it does not demonstrably reverse against SVM-RBF at
any size tested, it becomes parity rather than a lead against XGBoost on the raw
$122$ features, and whether rare-attack enrichment removes it is not decided by
this design, whose enriched arm ends before the reversal begins. A
lexicographic selection rule that charges for two-qubit gates returns
$n^{\ast}=4$ on NSL-KDD and $n^{\ast}=6$ on UNSW-NB15 without any change of
threshold, so the embedding dimension is a property of the dataset rather than a
constant to be inherited. And the operating window is closed from above: at
$K=80$, where the rule requires eight qubits, every one of $48$ comparisons
favours the classical model, with off-diagonal Gram dispersion decaying faster
for the \ZZ{} map than for its entanglement-free control, by the factor of two
that phase-term counting predicts at $K=20$ and by $3.3$ at $K=80$, and tracking
macro-$F_1$ at $r=+0.77$ over seven widths (95\% CI $[+0.04,+0.96]$).

That last result is the one we would most like others to test. It ties a
structural property of the circuit ($\binom{n}{2}$ pairwise phase terms
against $n$ single-qubit ones) to a decay exponent we can measure and to an
end-task metric that follows it down. If it holds on other corpora and other
maps, then the practical question for quantum kernels on tabular data is not how
to make them more expressive but how to keep them from concentrating, and the
useful width is bounded well below where hardware limits bite.

We are equally clear about what this study does not establish. The circuits were
not executed on a quantum processor, and ``NISQ-aware'' here means a circuit
budget chosen under an explicit cost model, not hardware validation. The
comparison covers two datasets, two tree ensembles and three kernel baselines,
not the deep tabular architectures. The favourable regime is narrow on every
axis we varied. And five claims from the earlier version of this work do not
survive the revised protocol; they are withdrawn in Sec.~\ref{sec:intro} and
corrected in Sec.~\ref{sec:erratum}.

Three directions follow directly. Executing the same kernels on hardware would
close the one gap our simulations cannot: the released code contains the
execution path and the three-tier comparison it would produce is already
specified. Testing the dispersion--accuracy relationship on projected quantum
kernels, which concentrate differently, would show whether the boundary we
measured is a property of fidelity kernels specifically or of the approach.
And since the crossover reproduces on one dataset and not the other, the
mechanism behind it, why more data helps this kernel more than it helps a
tree ensemble and only under a skewed prior, remains open. Our full
implementation, protocol files, intermediate artifacts and two independent audit
scripts are released so that each of these can be started from where we stopped.
\end{revbody}
